From the information shown in the results, it can be concluded that the local linear modeling approach is effective for system identification of the nonlinear CSTR. The step test data provided sufficient information to capture the dynamics of the system, and the resulting models demonstrated good predictive capabilities.

The main findings of this study can be summarized as follows:

\begin{itemize}
    \item The local linear models were able to accurately represent the nonlinear behavior of the CSTR within the operating range of interest.
    \item The use of industrial-style step test data proved to be a practical and efficient method for system identification in this context.
    \item The identified models can be utilized for advanced process control (APC) applications, providing a foundation for further optimization and control strategies.
    \item Future work could explore the integration of these models into real-time control systems and investigate their performance under varying operating conditions.
    \item In order to correct for the nonlinearities of the system, it is recommended to implement a gain scheduling approach or a nonlinear control strategy that can adapt to the changing dynamics of the CSTR.
\end{itemize}